\documentclass[review,3p,times]{elsarticle}

\usepackage{lineno}
\usepackage{amsmath,amssymb}
\usepackage{graphicx}
\usepackage{booktabs}
\usepackage{hyperref}
\usepackage{xurl}

\graphicspath{{figures/}}

\journal{MethodsX}

\begin{document}

\begin{frontmatter}

\title{iSAGE-CD: sparse point supervision for bi-temporal change detection
       with pair-anchored annotations}

\author[1]{Osmar Luiz Ferreira de Carvalho\corref{cor1}}
\ead{osmaracj@gmail.com}
\author[1]{Osmar Ab\'ilio de Carvalho J\'unior}
\author[1]{Anesmar Olino de Albuquerque}
\author[1]{Daniel Guerreiro e Silva}
\cortext[cor1]{Corresponding author}
\affiliation[1]{organization={University of Bras\'ilia (UnB)},
                city={Bras\'ilia},
                country={Brazil}}

\begin{abstract}
% [~300 words placeholder]
Change detection in remote sensing requires per-pixel labels on two
co-registered dates, a cost that scales poorly with area and event
frequency. iSAGE has shown that, in single-image semantic segmentation,
expert clicks targeting confident model errors recover dense supervision
quality at a fraction of a percent of the pixel budget. We extend that
philosophy to the bi-temporal case with iSAGE-CD: a session format,
annotation UI, and iteration protocol whose fundamental unit is the
\emph{pair} rather than the image. Labels are anchored to a
(reference, frame) pair, two synchronized canvases enforce a shared pixel
grid, and the annotation tree \emph{is} the dataset -- no separate export.
We describe the method and apply it to Amazon deforestation using Planet
imagery aligned to the PRODES 2022 calendar. Code is released under CC BY 4.0
at \url{https://github.com/osmarluiz/isage_cd}.
\end{abstract}

\begin{keyword}
change detection \sep sparse supervision \sep human-in-the-loop \sep
remote sensing \sep deforestation \sep annotation tool
\end{keyword}

\end{frontmatter}

% ------------------------------------------------------------------
% Specifications table (MethodsX required)
% ------------------------------------------------------------------
% Specifications table (MethodsX required)
% To be finalized. Layout follows the MethodsX template.
\begin{table}[!h]
\centering
\begin{tabular}{p{0.30\linewidth}p{0.60\linewidth}}
\toprule
Subject area              & Environmental science, Remote sensing \\
More specific subject     & Change detection; deep learning; annotation tools \\
Method name               & \textbf{iSAGE-CD} \\
Name/reference of original method & iSAGE \cite{isage} extended to bi-temporal
                            change detection \\
Resource availability     & Code: \url{https://github.com/osmarluiz/isage_cd}
                            (CC BY 4.0) \\
Related research article  & None. \\
\bottomrule
\end{tabular}
\end{table}


% ------------------------------------------------------------------
% Background (max 500 words)
% ------------------------------------------------------------------
\section{Background}
% =============================================================
% BACKGROUND — max 500 words (MethodsX requirement)
% =============================================================

Change detection is one of the most widely used operations in remote
sensing. Comparing two co-registered images of the same place at
different dates reveals where the surface has changed, and the same
operation serves many fields. It maps urban
expansion~\cite{chen2020spatial}, assesses damage after
disasters~\cite{gupta2019xbd}, tracks agricultural and land-cover
transitions~\cite{yang2021asymmetric}, and monitors
deforestation~\cite{prodes,mapbiomas}. The deep-learning methods behind
these applications~\cite{daudt2018fully}, and their shared challenges,
are the subject of dedicated surveys~\cite{khelifi2020deep,shi2020change}.

Supervised change detection is built on pixel-level labels, and
producing them is the main cost of a dataset. An analyst inspects the
imagery, decides the state of each pixel, and paints a dense mask.
Change makes this harder than single-image labeling in two ways. The
unit is a pair of co-registered dates rather than one image, so the
analyst reasons about both at once, and the same scene can be labeled
against several reference dates, which multiplies the work. Expert time
becomes the binding constraint, and reducing it has driven a range of
label-efficient strategies. Weakly supervised methods replace dense
masks with image-level tags or
scribbles~\cite{lin2016scribblesup,hua2021semantic}. Semi-supervised
methods combine a small labeled set with a large unlabeled
one~\cite{zhou2026decaf}. Active learning chooses which samples the
human should label next~\cite{ren2021survey}. Human-in-the-loop
annotation lets the model and the user take turns, the model proposing a
mask and the user correcting it~\cite{amershi2014power,lenczner2022dial}.
Point supervision replaces dense masks with a small set of
clicks~\cite{Bearman2016}, and recent work brings it to change detection
specifically~\cite{fang2024point}.

iSAGE~\cite{isage} (\emph{Iterative Sparse Annotation Guided by Expert})
is a recent point-based, human-in-the-loop framework for remote-sensing
semantic segmentation. An expert clicks only on the model's confident
mistakes and an error-weighted loss amplifies the gradient at each
click, with no pseudo-labels, conditional random fields, or
foundation-model prompts, and the annotation record is itself the
dataset. On single-image benchmarks it recovered most of the quality of
dense labeling from far less than 1\% of the pixels. iSAGE was designed
for the single-image case, and its label file, click loop, and
error-weighted loss all assume one image per annotation. This article
introduces \textbf{iSAGE-CD}, an adaptation of iSAGE to change
detection. iSAGE-CD is a general annotation tool that applies to any
bi-temporal change task on any co-registered imagery. We describe the
tool, its pair-anchored session format, and the annotation workflow, and
we illustrate them with a working annotation session. The aim is to
document a reusable method rather than to run a controlled experiment,
and a full quantitative evaluation is the subject of a separate
manuscript in preparation.


% ------------------------------------------------------------------
% Method details
% ------------------------------------------------------------------
\section{Method details}
% TODO: draft 04_method_details


% ------------------------------------------------------------------
% Method validation
% ------------------------------------------------------------------
\section{Method validation}
% TODO: draft 05_method_validation


% ------------------------------------------------------------------
% Limitations
% ------------------------------------------------------------------
\section{Limitations}
% =============================================================
% LIMITATIONS
% =============================================================

The validation reported in this article is qualitative and limited
to a single scene and a single event (the PRODES 2022 window over
a subset of the Brazilian arc of deforestation). All annotations
were placed by a single expert; no inter-annotator agreement, no
comparison against dense supervision, and no held-out test set
are provided. A quantitative evaluation with multiple annotators
and against dense ground truth is the subject of a separate
manuscript in preparation. Nothing in this article should be read
as a scientific claim about deforestation detection accuracy; the
validation establishes that the tool and the protocol produce
usable output, not that the output is optimal.

The tool assumes that the two images of a pair share the same
pixel grid. Co-registration and resampling are upstream concerns
and are not addressed by iSAGE-CD. Similarly, only bi-temporal
change (one reference and one frame per pair) is supported;
multi-date time series would require an extension of the
pair-anchored label format described in
Section~\ref{sec:label_format}.

The interface is a local PyQt5 desktop application and was tested
primarily on Windows. Training and prediction scripts are supplied
separately from the annotator itself; users adopting iSAGE-CD for
a different task must plug in their own training loop and their
own inference script, both configured to produce the per-frame
prediction PNGs in the class-index format described in
Section~\ref{sec:iteration}.

Finally, labels in iSAGE-CD are anchored to a pair rather than to
an image, which means the same visual polygon may carry different
labels under different references (for example, a cleared area
may be \emph{new} against one before-image and \emph{old} against
another). This is a design feature for change detection but is a
semantic shift relative to conventional single-image annotation
and should be accounted for when training analysts on the tool.


% ------------------------------------------------------------------
% Ethics / declarations / acknowledgements
% ------------------------------------------------------------------
\section{Ethics statements}
% =============================================================
% ETHICS STATEMENTS
% =============================================================

The work reported in this article involved no human subjects and
no animal experiments; the annotation session described in
Section~\ref{sec:val_session} was carried out by one of the
authors, so no ethics review board approval was required.

The licensing of the imagery, of the annotations and of the tool
itself is stated under Data availability.


\section{Declaration of competing interest}
The authors declare no known competing financial interests or personal
relationships that could have influenced the work reported in this
article.

\section{Acknowledgements}
% =============================================================
% ACKNOWLEDGEMENTS
% =============================================================

The authors are grateful for the financial support of the grant
from the National Council for Scientific and Technological
Development (CNPq) for Professor Osmar Ab\'ilio de Carvalho
J\'unior. This study was partly financed by the Coordination for
the Improvement of Higher Education Personnel (CAPES) (Finance
Code~001), the Research Support Foundation of the Federal
District (FAPDF) (Project 00193.00002237/2022-92), and CNPq
(research projects 312608/2021-7, 421413/2023-9, and
421291/2022-2).


% ------------------------------------------------------------------
% References
% ------------------------------------------------------------------
\bibliographystyle{elsarticle-num}
\bibliography{references}

\end{document}
